%%
%% thesis.tex
\documentclass[newenv,final]{ncuthesis}
%%
\ncusetup {
  info = {
    CLC             = {O}, % 可省略花括号
    secret-level*   = 公开,
    UDC             = 123.123.1,
    studentid       = 402200220047,
    degree          = 硕士,
    title           = 南昌大学学位论文模板使用演示,
    title*          = Demonstration\\of\\Nanchang University Thesis Template,
    author          = 朱金宝,
    department      = 物理与材料学院,
    supervisor      = 舒富文教授,
    field           = 理学,
    major           = 天体物理,
    defence-date    = 2023 年 2 月 16 日,
    tdc-chairman    = 主席,
    reviewer        = {评阅人 1, 评阅人 2},
    completion-date = 2023 年 2 月 16 日
  }
}
%%
%%
\begin{document}
\maketitle
%%
\frontmatter
\makedecaut
%%
\begin{abstract}
这是摘要内容。
\keywords{关键词 1, 关键词 2, 关键词 3}
\end{abstract}
\begin{abstract*}
This is abstract.
\keywords{kw1, kw2, kw3}
\end{abstract*}
%%
\tableofcontents
%%
\mainmatter
\chapter{引言}
\section{历史}
我们知道 lshort\cite{lshortcn} 是入门 \LaTeX\ 的必读文献。\footnote{test}

\begin{figure}[htb]
\centering
\begin{subfigure}{.45\textwidth}
\centering
\includegraphics[width=5cm]{figures/test.pdf}
\subcaption{这是子图}
\label{fig:subfigexp-a}
\end{subfigure}
\begin{subfigure}{.45\textwidth}
\centering
\includegraphics[width=5cm]{figures/test.pdf}
\subcaption{这又是子图}
\label{fig:subfigexp-b}
\end{subfigure}
\begin{subfigure}{.45\textwidth}
\centering
\includegraphics[width=5cm]{figures/test.pdf}
\subcaption{这还是子图}
\label{fig:subfigexp-c}
\end{subfigure}
\caption{子图排版示例}
\label{fig:subfig}
\end{figure}

\begin{law}[黑洞面积定律]
  黑洞事件视界的面积不可能随着时间的推移而减小。
\end{law}
\begin{proof}
即得易见平凡，仿照上例显然。留作习题答案略，读者自证不难。反之亦然同理，推论自然成立。略去过程 QED，由上可知证毕。
\end{proof}
\appendix
\chapter{公式推导}
\section{Einstein 场方程}
%% 推导过程
%%
\backmatter
\bibliography{references.bib}
%%
\begin{acknowledgements}
感谢……
\signoff{朱金宝}{2022 年 12 月 31 日}
\end{acknowledgements}
%%
%%
%% 注意，以下内容本科生论文应当删除，否则报错！
%%
\begin{researchresults}
\begin{published}
\item 条目
\end{published}
\begin{tobepublished}
\item 条目
\end{tobepublished}
\begin{reports}
\item 条目
\end{reports}
\begin{others}[另外还有] % 标题作为可选参数
\item 条目
\end{others}
\begin{papers}
\item 条目
\end{papers}
\end{researchresults}
\end{document}
%% 结束
%%
